%% =====================================================================
%%  Bilan 5 — Hier et aujourd'hui (unités 17 à 20)
%%  Pistes : B5.01 (écoute notée) · B5.02 (expression orale)
%% =====================================================================
\bilan{5}

\noindent{\fontsize{11}{14.6}\selectfont Dernière partie terminée : bravo ! Vous savez maintenant
raconter votre week-end ou un voyage, dire où vous étiez et parler de vous au passé, au présent et au
futur. Révisez l'essentiel, puis faites les huit exercices notés sur 20. Les corrigés sont
page~\pageref{cor-b5-1}.\par}

%% ---------------------------------------------------------------------
\rubrique{grammaire}{Révisez}

\begin{paire}
\begin{regle}[\en{was} et \en{were} · unité 17]
\en{I / he / she / it was} · \en{you / we / they were}\par
\en{wasn't} · \en{weren't} · \en{Were you…? — Yes, I was.}\par\vspace{1mm}
\en{there was} / \en{there were} · \en{two years ago} · \en{I was born in…}
\end{regle}
\begin{regle}[Le prétérit · unités 18 et 19]
Réguliers : \en{watched}, \en{stayed}, \en{visited}, \en{tried}, \en{stopped}\par
Irréguliers : \en{went}, \en{saw}, \en{took}, \en{ate}, \en{bought}, \en{had}\par\vspace{1mm}
\en{I didn't go} · \en{Did you go? — Yes, I did.}\par\vspace{1mm}
{\fontsize{9.6}{12}\selectfont Après \en{did} et \en{didn't} : la base verbale.\par}
\end{regle}
\end{paire}

\begin{regle}[Les temps et les questions · unité 20]
\en{I work} (habitude) · \en{I'm working} (maintenant) · \en{I'm going to work} (projet) · \en{I worked} (passé)\par\vspace{1mm}
Questions : auxiliaire + sujet : \en{Do…?} · \en{Are…?} · \en{Can…?} · \en{Was…?} · \en{Did…?}
\end{regle}

%% ---------------------------------------------------------------------
\rubrique{pratique}{Testez-vous}

\exercice[ecrire,niveau=1,points=2]{\en{was} ou \en{were} ?}{Complétez. \textit{(0,5 point par
réponse)}}
\begin{anglais}
\begin{questions}[2]
\item I \trou[14mm]{was} tired yesterday.
\item My friends \trou[14mm]{were} at the party.
\item \trou[14mm]{Were} you at home last night?
\item It \trou[14mm]{wasn't} very cold.
\end{questions}
\end{anglais}

\exercice[ecrire,niveau=1,points=2]{Le prétérit}{Écrivez le prétérit de chaque verbe.
\textit{(0,5 point par réponse)}}
\begin{anglais}
\begin{questions}[4]
\item watch → \trou[18mm]{watched}
\item study → \trou[18mm]{studied}
\item go → \trou[18mm]{went}
\item buy → \trou[18mm]{bought}
\end{questions}
\end{anglais}

\exercice[ecrire,niveau=2,points=3]{Négations et questions}{Transformez chaque phrase : (?) en question,
(−) en phrase négative. \textit{(0,5 point par réponse)}}
\begin{questions}
\item We went to the beach. (−) \hfill \enb{\trou[66mm]{We didn't go to the beach.}}
\item She called you. (?) \hfill \enb{\trou[66mm]{Did she call you?}}
\item I saw the movie. (−) \hfill \enb{\trou[66mm]{I didn't see the movie.}}
\item They were happy. (?) \hfill \enb{\trou[66mm]{Were they happy?}}
\item He was at home. (−) \hfill \enb{\trou[66mm]{He wasn't at home.}}
\item You took the bus. (?) \hfill \enb{\trou[66mm]{Did you take the bus?}}
\end{questions}

\exercice[ecrire,niveau=2,points=3]{Le bon temps}{Mettez le verbe au temps qui convient.
\textit{(0,5 point par réponse)}}
\begin{questions}
\item Last year, we \trou[40mm]{visited} London. \hfill (\en{visit})
\item Right now, Léa \trou[40mm]{is reading}. \hfill (\en{read})
\item Every Sunday, I \trou[40mm]{have} brunch. \hfill (\en{have})
\item Tomorrow, they \trou[40mm]{are going to fly} to Boston. \hfill (\en{fly})
\item Two days ago, Sam \trou[40mm]{made} a cake. \hfill (\en{make})
\item Yesterday, it \trou[40mm]{was} cold. \hfill (\en{be})
\end{questions}

\exercice[ecouter,niveau=2,points=3,piste=B5.01]{Le week-end de Sam}{Sam raconte son week-end à
Atlanta. Écoutez, puis indiquez si chaque phrase est vraie (V) ou fausse (F).
\textit{(0,5 point par réponse)}}
\begin{anglais}
\begin{questions}[2]
\item Sam and Maya went to Atlanta by plane. \hfill\trou[9mm]{V}
\item They visited their grandparents. \hfill\trou[9mm]{F}
\item Their mother cooked on Saturday. \hfill\trou[9mm]{V}
\item On Sunday, they saw a basketball game. \hfill\trou[9mm]{F}
\item The Braves won the game. \hfill\trou[9mm]{F}
\item They came back on Monday. \hfill\trou[9mm]{V}
\end{questions}
\end{anglais}

\exercice[lire,niveau=2,points=3]{Le journal de Léa}{Lisez la page du journal de Léa, puis indiquez si
chaque phrase est vraie (V) ou fausse (F). \textit{(0,5 point par réponse)}}
\begin{anglais}
\begin{tcolorbox}[otBase,colback=otOrangeTint!40,colframe=otOrangeTint!40,fontupper=\fontsize{10.8}{14.4}\selectfont\itshape]
{\otSemi\upshape\fontsize{9.4}{11}\selectfont\color{otOrangeDark}March 16th}\par\vspace{1mm}
What a day! Yesterday was my birthday. Emma asked me to meet her at Sam's café at seven. When I opened
the door, everybody was there: Jack, Sam, Maya, Mrs.\ Lopez… and even Mr.\ Carter! Sam made a
chocolate cake, and Jack gave me a book about Route 66. We danced until midnight. Today, I'm very
tired… but so happy!
\end{tcolorbox}
\begin{questions}[2]
\item Léa's birthday was on March 15th. \hfill\trou[9mm]{V}
\item The party was at the studio. \hfill\trou[9mm]{F}
\item Mr.\ Carter was at the party. \hfill\trou[9mm]{V}
\item Jack made the cake. \hfill\trou[9mm]{F}
\item They danced until midnight. \hfill\trou[9mm]{V}
\item Today, Léa isn't tired. \hfill\trou[9mm]{F}
\end{questions}
\end{anglais}

\exercice[ecrire,niveau=2,points=2]{Posez la question}{Écrivez la question qui correspond à chaque
réponse. \textit{(0,5 point par réponse)}}
\begin{questions}
\item I was born in 1998. \hfill \enb{\trou[66mm]{When were you born?}}
\item We went to Boston. \hfill \enb{\trou[66mm]{Where did you go?}}
\item I'm reading a book. \hfill \enb{\trou[66mm]{What are you doing?}}
\item I'm going to travel. \hfill \enb{\trou[66mm]{What are you going to do?}}
\end{questions}

\exercice[ecrire,niveau=3,points=2]{Votre année}{En quatre phrases, racontez une chose que vous avez
faite l'année dernière, décrivez une habitude et présentez un projet. \textit{(0,5 point par phrase
correcte)}}
\lignes{4}

%% ---------------------------------------------------------------------
\rubrique{ecouter}{Parlez}

\exercice[ecouter,parler,niveau=3,piste=B5.02]{Le grand entretien}{Pour clore ce premier tome,
Mr.\ Carter vous interroge comme pour un entretien. Répondez \textbf{à voix haute} pendant chaque pause.
Cet exercice n'est pas noté : comparez vos réponses avec les modèles du corrigé.}
\begin{center}\otTableau
\begin{tabular}{@{}>{\otSemi\fontsize{8.6}{10}\selectfont\color{otNavy}}l@{\hspace{3mm}}>{\enOn}p{70mm}<{\enOff}@{\hspace{6mm}}p{78mm}@{}}
 & \otTH{Mr.\ Carter dit\ldots} & \otTH{Vous\ldots}\\\midrule
1. & Hello! Can you introduce yourself? & donnez votre nom, votre âge et votre ville.\\
2. & Where were you born? & répondez.\\
3. & What did you do last year? & racontez un événement.\\
4. & What do you do in your free time? & parlez de vos loisirs.\\
5. & And what are you going to do next year? & présentez un projet.\\\bottomrule
\end{tabular}\end{center}

\otScore{20}
